%% LyX 2.4.4 created this file.  For more info, see https://www.lyx.org/.
%% Do not edit unless you really know what you are doing.
\documentclass[english]{article}
\usepackage[T1]{fontenc}
\usepackage[latin9]{inputenc}
\usepackage{amsmath}
\usepackage{amsthm}
\usepackage{amssymb}
\usepackage{geometry}
\geometry{verbose,tmargin=1in,bmargin=1in,lmargin=1in,rmargin=1in}
\PassOptionsToPackage{normalem}{ulem}
\usepackage{ulem}

\makeatletter
%%%%%%%%%%%%%%%%%%%%%%%%%%%%%% Textclass specific LaTeX commands.
\numberwithin{equation}{section}
\numberwithin{figure}{section}
\newlength{\lyxlabelwidth}      % auxiliary length 
\theoremstyle{plain}
\newtheorem{thm}{\protect\theoremname}[section]
\theoremstyle{definition}
\newtheorem{xca}[thm]{\protect\exercisename}
\theoremstyle{remark}
\newtheorem*{rem*}{\protect\remarkname}

%%%%%%%%%%%%%%%%%%%%%%%%%%%%%% User specified LaTeX commands.
\usepackage{fancyhdr}
\usepackage{lastpage}
\pagestyle{fancy}
\usepackage{enumitem}
\usepackage{ifthen}
\everymath=\expandafter{\the\everymath\displaystyle}
%\DeclareMathSizes{10}{10}{10}{10}
\usepackage{multicol}

\usepackage{tikz}
\usetikzlibrary{patterns}
\usetikzlibrary{plotmarks}
\usepackage{pgfplots}
\pgfplotsset{compat=newest}

\usepackage{calc}
\usepackage[nomessages]{fp}

\usepackage{datetime}
% Added by lyx2lyx
\setlength{\parskip}{\smallskipamount}
\setlength{\parindent}{0pt}

\makeatother

\usepackage{babel}
\providecommand{\exercisename}{Exercise}
\providecommand{\remarkname}{Remark}
\providecommand{\theoremname}{Theorem}

\begin{document}
\renewcommand{\author}{Dr.\,James Ashe}

\newcommand{\classNum}{439}

\newcommand{\classSec}{A}

\newcommand{\nameType}{Name: \hspace{-4cm}}

\newcommand{\examType}{Homework 3}

\newcommand{\Date}{\formatdate{6}{2}{2014}} %%\AdvanceDate[12] \today

%\newdate{my_date}{\day}{\month}{\year}%   
%\AdvanceDate[-5] 
%\displaydate{my_date}

%%First page's header%%%%%%%%%%%%%%%%%%%%%%
\fancypagestyle{firststyle} %%header settings for first pagr
{
	\fancyhead[L]{MTH \classNum{}(\classSec{}) \author{}\\
	\textbf{\examType{}} \Date{}}
	\fancyhead[C]{\textbf{\nameType}}
	\setlength{\headsep}{14pt}
}
%%%%%%%%%%%%%%%%%%%%%%%%%%%%%%%%%%%%%%%%%%

%%Next page's header%%%%%%%%%%%%%%%%%%%%%%%
%\setlist{nolistsep} %eliminates space between list items
\lhead{MTH \classNum{}(\classSec{})} 
\chead{\textbf{\nameType}} 
\cfoot{\thepage{}~of~\pageref{LastPage}} 
\setlength{\headsep}{5pt} 
%%%%%%%%%%%%%%%%%%%%%%%%%%%%%%%%%%%%%%%%%%%

%%Various settings; see the preamble for other settings%%%%%
%%\setenumerate[0]{leftmargin=12pt,itemindent=0pt,labelwidth=0pt}
\renewcommand{\labelenumi}{\textbf{\arabic{enumi}.}}
\renewcommand{\labelenumii}{\textbf{\alph{enumii}.}}
\thispagestyle{firststyle}
%%%%%%%%%%%%%%%%%%%%%%%%%%%%%%%%%%%%%%%%%%

\null

\vspace{1cm}

\setcounter{section}{3}

\Large{\textbf{Chapter 3}} \large

All excercises \uline{without }a $\bigstar$ from 3.1--3.13

\textbf{\textbf{\examType} due: \formatdate{18}{3}{2014}}
\begin{xca}
Prove that $0_{R}$ and $R$ are ideals of $R$ for any ring $R$.
The ideal $0_{R}$ is called the \emph{zero ideal}. 

\begin{proof}
We already know that $0_{R}$ and $R$ are non-empty subrings of $R$.
The only element in $0_{R}$ is $0_{R}$, and $a0_{R}=0_{R}a=0_{R}$
for all $a\in R$. So $0_{R}$ is an ideal. Also, $ar$ and $ra$
are both in $R$ for any $r,a\in R$. So $R$ is an ideal as well.
\end{proof}
\end{xca}

\begin{xca}
Let $\mathbb{R}[x]$, the set of all polynomials with real coefficients.
Let $C$ be the subset of $\mathbb{R}[x]$ with constant terms that
are $0$. For example, $x^{3}+x\in C$ but $x^{3}+x-4\notin C$. Prove
that $C$ is or is not an ideal.

\begin{proof}
$C$ is an ideal. $x\in C$ so then $C$ is non-empty. 

Let $i,j\in C$. $i-j\in C$ since both $i$ and $j$ have a constant
term of $0$, and so the constant term of $i-j$ is $0$. 

Any polynomial in $\mathbb{R}[x]$ can be written as $f(x)+c$ where
$c$ is the constant term and $f(x)\in C$. Let $f(x)+0\in C$ and
$g(x)+c\in\mathbb{R}[x]$. $f(x)(g(x)+c)=f(x)g(x)+f(x)c$ which has
a constant term of $0$. So then $C$ is an ideal by the ideal test.
\end{proof}
\end{xca}

\begin{xca}
Prove that the set $I=\{(a,0)\colon a\in\mathbb{Z}\}$ is an ideal
of the ring $\mathbb{Z}\times\mathbb{Z}$. 

\begin{proof}
$\mathbb{Z}\times\mathbb{Z}\neq\emptyset$ since $(1,0)\in\mathbb{Z}\times\mathbb{Z}$.
Let $(a,0),(b,0)\in\mathbb{Z}\times\mathbb{Z}$. $(a,0)-(b,0)=(a-b,0)\in\mathbb{Z}\times\mathbb{Z}$.
Now let $(x,y)\in\mathbb{Z}\times\mathbb{Z}$. $(x,y)(a,0)=(x\cdot a,y\cdot0)=(x\cdot a,0)\in\mathbb{Z}\times\mathbb{Z}$.
So $\mathbb{Z}\times\mathbb{Z}$ is an ideal by the ideal test.
\end{proof}
\end{xca}

\begin{xca}
Let $A$ and $B$ be ideals of $R$. Prove that $A\cap B$ is also
an ideal of $R$.

\begin{proof}
$A$ and $B$ are both ideals. Thus $0_{R}\in A$ and $0_{R}\in B$
and $0_{R}\in A\cap B$ ($0_{R}\in I$ for any ideal in $R$. If not
then $a\cdot0_{R}=0_{R}\notin I$ for any $a\in I$.) 

Let $x,y\in A\cap B$. $x-y\in A$ and $x-y\in B$ since $A$ and
$B$ are both subrings, and subrings are closed under subtraction. 

Let $r\in R$. $rx,xr\in A$ and $ry,yr\in B$ since $A$ and $B$
are ideals. So each of these products are also in $A\cap B$, and
$A\cap B$ is an ideal by the ideal test.
\end{proof}
\end{xca}

\begin{xca}
Let $A$ and $B$ be ideals of $R$. Prove that $A\cup B$ is not
necessarily an ideal of $R$.

\begin{proof}
Consider the ring $\mathbb{Z}$ with ideals $2\mathbb{Z}$ and $3\mathbb{Z}$.
$2\in2\mathbb{Z}$ and $3\in3\mathbb{Z}$, but $3-2=1\notin2\mathbb{Z}\cup3\mathbb{Z}$.

\begin{rem*}
The union of plenty ideals \emph{is }an ideal, e.g., $I\cup0_{R}$
and $I\cup I$ for any ideal $I$ in $R$. Recall that the union of
subrings is not necessarily a subring.
\end{rem*}
\end{proof}
\end{xca}

\begin{xca}
If $I$ is an ideal in the ring $R$ and $J$ is an ideal of the ring
$S$, prove that $I\times J$ is an ideal of the ring $R\times S$.

\begin{proof}
$I,J\neq\emptyset$ as they are both ideals. So there exists and $i\in I$
and $j\in J$ such that $(i,j)\in I\times J$, and $I\times J\neq\emptyset$.

Let $i,i'\in I$ and $j,j'\in J$. $(i,j)-(i',j')=(i-i',j-j')\in I\times J$
as $i-i'\in I$ and $j-j'\in J$ since $I$ and $J$ are both ideals.

Let $(r,s)\in R\times S$. $(i,j)(r,s)=(ir,js)\in R\times S$. Similarly
$(r,s)(i,j)\in R\times S$. So $I\times J$ is an ideal by the ideal
test. 
\end{proof}
\end{xca}

\begin{xca}
Let $R$ be a ring with unity $1_{R}$ and let $I$ be an ideal of
$R$. If $1_{R}\in I$ prove $I=R$.  

\begin{proof}
$I\subset R$ by definition of an ideal. If $1_{R}\in I$ then $1_{R}r,\,r1_{R}\in I$
since $I$ is an ideal. So then $1_{R}r=r1_{R}=r\in I$ for all $r\in R$
which implies that $R\subset I$ and $R=I$.
\end{proof}
\end{xca}

\begin{xca}
Let $R$ be a ring with unity $1_{R}$ and let $I$ be an ideal of
$R$. If $I$ has an element $x$ and $x^{-1}$ such that $x\cdot x^{-1}=x^{-1}\cdot x=1_{R}$,
prove that $I=R$.

\begin{proof}
So if $I$ has two elements such that $x\cdot x^{-1}=x^{-1}\cdot x=1_{R}$
then $1_{R}\in I$ since $I$ is a subring and closed under multiplication.
It follows that $I=R$ by the exercise above.
\end{proof}
\end{xca}

\begin{xca}
$\bigstar$
\end{xca}

\begin{xca}
$\bigstar$
\end{xca}

\begin{xca}
$\bigstar$
\end{xca}

\begin{xca}
Show that $5\mathbb{Z}$ is a prime ideal. 

\begin{proof}
We know from example 3.1 that $5\mathbb{Z}$ is an ideal in $\mathbb{Z}$.
$6\notin5\mathbb{Z}$, so it is a proper ideal. Let $a,b\in\mathbb{Z}$
such that $ab\in5\mathbb{Z}$. So then $ab=5n$ for some $n\in\mathbb{Z}$
which implies that $5|ab$. Suppose that $b\notin5\mathbb{Z}$ so
that $5\nmid b$. Since $5$ is prime, $5|a$ (by Euclid's Lemma).
Thus $5\mathbb{Z}$ is a prime ideal.
\end{proof}
\end{xca}


\end{document}
